\SixthArticle{17}
Die Untersuchung der Beschaffenheit der Summe der Reihe $F(\alpha,\beta,\gamma,1)$ beschränkt sich daher ihrer Natur nach auf den Fall, dass $\gamma-\alpha-\beta$ eine positive Größe ist; dann stellt jene Summe stets eine endliche Größe dar. Wir schicken jedoch folgende Bemerkung voraus. Wenn die Koeffizienten der Reihe $1+ax+bxx+cx^3+$ usw. $=S$ von einem gewissen Glied an gegen Null abnehmen, so muss das Produkt
\[
(1-x)S=1+(a-1)x+(b-a)xx+(c-b)x^3+\quad\text{usw.}
\]
für $x=1$ als $=0$ angesetzt werden, selbst wenn die Summe der Reihe $S$ selbst unendlich groß wird. Denn fasst man zwei Glieder zusammen, so ist die Summe $=a$, bei drei Gliedern $=b$, bei vier Gliedern $=c$ usw.; die Grenze der ins Unendliche fortgesetzten Summe ist daher $=0$. Ist also $\gamma-\alpha-\beta$ eine positive Größe, so muss $(1-x)F(\alpha,\beta,\gamma-1,x)=0$ für $x=1$ gesetzt werden, woraus wir nach Gleichung 15 des Artikels 7 erhalten:
\[
0=\gamma(\alpha+\beta-\gamma)F(\alpha,\beta,\gamma,1)+(\gamma-\alpha)(\gamma-\beta)F(\alpha,\beta,\gamma+1,1),\quad\text{oder}
\]
\[
[37]\qquad F(\alpha,\beta,\gamma,1)=\frac{(\gamma-\alpha)(\gamma-\beta)}{\gamma(\gamma-\alpha-\beta)}F(\alpha,\beta,\gamma+1,1)
\]
Da sich in gleicher Weise
\[
\begin{aligned}
F(\alpha,\beta,\gamma+1,1)&=\frac{(\gamma+1-\alpha)(\gamma+1-\beta)}{(\gamma+1)(\gamma+1-\alpha-\beta)}F(\alpha,\beta,\gamma+2,1)\\[6pt]
F(\alpha,\beta,\gamma+2,1)&=\frac{(\gamma+2-\alpha)(\gamma+2-\beta)}{(\gamma+2)(\gamma+2-\alpha-\beta)}F(\alpha,\beta,\gamma+3,1)
\end{aligned}
\]
ergibt und so fort, gilt allgemein, wenn $k$ eine beliebige positive ganze Zahl bezeichnet,
\SourceLeaf{153}
\[
\begin{array}{@{}ll@{}}
&F(\alpha.\beta,\gamma,1)\quad\text{ist gleich dem Produkt aus}\quad F(\alpha,\beta,\gamma+k,1)\\[6pt]
\text{und}&(\gamma-\alpha)(\gamma+1-\alpha)(\gamma+2-\alpha)\mathinner{.\,.\,.\,.}(\gamma+k-1-\alpha)\\[5pt]
\text{und}&(\gamma-\beta)(\gamma+1-\beta)(\gamma+2-\beta)\mathinner{.\,.\,.\,.}(\gamma+k-1-\beta)\\[8pt]
\multicolumn{2}{@{}l@{}}{\text{geteilt durch das Produkt}}\\[6pt]
\text{aus}&\gamma(\gamma+1)(\gamma+2)\mathinner{.\,.\,.\,.}(\gamma+k-1)\\[5pt]
\text{und}&(\gamma-\alpha-\beta)(\gamma+1-\alpha-\beta)(\gamma+2-\alpha-\beta)\mathinner{.\,.\,.}(\gamma+k-1-\alpha-\beta)
\end{array}
\]
